% ================================================
% 文件: 04_电视接收机.tex
% 章节: 第9章 电视接收机
% 所属部分: 接收机技术
% 版本: v1.0
% ================================================
% 本章目录结构（树状）:
% \chapter{电视接收机}
% ├── \section{电视接收机概述}
% ├── \section{模拟电视接收机}
% ├── \section{数字电视接收机}
% ├── \section{电视接收机电路原理}
% └── \section{电视接收机调试与维护}
% ================================================

\chapter{电视接收机}
\label{chap:television_receiver}

\section{电视接收机概述}
\label{sec:tv_receiver_intro}

电视接收机是一种能够接收并还原电视信号的设备。
它将射频电视信号转换为图像和声音输出。

电视接收机的发展历程：
\begin{itemize}
    \item \textbf{黑白电视}：1920年代-1950年代
    \item \textbf{彩色电视}：1950年代至今
    \item \textbf{数字电视}：2000年代至今
\end{itemize}

电视信号包含视频信号和音频信号，
通过调制方式传输。

电视接收机的根本任务，
是把空中传播的高频已调制电视信号还原为显像设备（早期为显像管，
现代为液晶、
等离子或 OLED 屏）所需的亮度、
色度信号，
以及扬声器所需的伴音。
与声音广播不同，
电视信号的信息量巨大：
除需传送与人眼视觉相对应的亮度（黑白）信息外，
还要传送色度（彩色）信息与行、
场同步信息，
因此对信道的带宽、
线性与信噪比要求远较声音广播严苛。
按照信号形式，
电视接收机的发展可划分为模拟电视与数字电视两大阶段，
二者在调制方式、
信号处理与整机结构上差异显著。

模拟电视时代以地面广播与有线电视（CATV）为主，
典型扫描体制为隔行扫描（如 625 行/50 场，
对应 PAL 与 SECAM；
525 行/60 场，
对应 NTSC）。
进入数字电视时代后，
信源由模拟视频转为经压缩编码的数字码流，
信道则采用正交频分复用（OFDM）或正交幅度调制（QAM）等抗多径、
高频谱效率的调制，
接收机内部大量工作转移到数字解调与解码芯片完成。
机顶盒（STB）的普及使“显示终端”与“接收解码单元”分离，
老旧模拟电视机也可借助机顶盒收看数字节目。

无论模拟还是数字，
电视接收机前端都遵循超外差基本思想：
将不同频道的射频信号统一变频到固定的中频，
再进行后续解调。
其变频关系仍为

\[
f_{\mathrm{IF}}=\left| f_{\mathrm{LO}}-f_{\mathrm{RF}} 
\right|
\]

只是电视中频数值、
带宽与伴音处理方式随制式而不同，
这正是后续各节讨论的重点。

\section{模拟电视接收机}
\label{sec:analog_tv_receiver}

模拟电视接收机接收模拟电视信号，
主要采用PAL、
NTSC或SECAM制式。

组成部分：
\begin{itemize}
    \item \textbf{调谐器}：选择电视频道
    \item \textbf{中频放大器}：放大中频信号
    \item \textbf{视频检波器}：提取视频信号
    \item \textbf{音频检波器}：提取音频信号
    \item \textbf{显像管/显示器}：显示图像
    \item \textbf{扬声器}：播放声音
\end{itemize}

模拟电视的缺点是易受干扰，
图像质量受信号强度影响较大。

模拟电视接收机的工作起点是高频调谐器（俗称“高频头”）。
它完成频道选择、
射频放大与第一次变频，
输出一个频率固定的图像中频（典型值随制式而变，
PAL-D/K 制常取 38.0MHz 量级，
NTSC-M 取 45.75MHz 量级）。
由于电视采用“残留边带”（vestigial sideband, VSB）调制，
视频信号的大部分下边带被抑制、
仅保留一小部分，
因而图像中频放大器的幅频特性必须按残留边带的特定形状修正，
才能保证亮度信号不发生线性失真。

视频检波通常采用同步检波或包络检波，
从图像中频中恢复出含亮度、
复合同步与色度信息的全电视信号（composite video）。
全电视信号经视频放大后一路送显像电路产生亮度，
另一路经钳位恢复直流电平，
并通过同步分离电路取出行、
场同步脉冲去控制扫描。
伴音则常与图像载波一同传输：
在接收端，
图像载频与伴音载频相差一个固定的“伴音内载波”频率（PAL-D/K 为 6.5MHz，
NTSC-M 为 4.5MHz），
由图像中频与伴音中频在检波级差拍产生第二伴音中频，
再经鉴频得到音频。

关于彩色，模拟彩色电视还需还原色度。以 PAL（逐行倒相）制为例，色度信号由两个正交分量 U（B-Y）与 V（R-Y）调制在约 4.43MHz 的彩色副载波上。PAL 的核心思想是让 V 分量的相位逐行倒相 180°（即第 n 行与第 n+1 行相位相反），在接收端再逐行反转回来；这样，因相位误差引起的色调失真在相邻两行平均后被抵消，从而获得稳定的彩色重现。色度副载波频率 $f_{\mathrm{sc}}$ 在 PAL 中通常取约 4.4336MHz，并需与行频保持严格的锁定关系（梳状滤波器利用行间相关性分离 U、V）。

模拟电视对噪声与多径极为敏感：
场强稍弱便出现“雪花”噪声，
多径反射则产生重影（ghosting），
强邻频干扰会引起网纹。
正因如此，
模拟电视的覆盖与质量高度依赖发射功率、
天线与传播环境。

模拟电视接收机的中频之后还需完成一系列模拟处理。
中放级通常受 AGC 控制：
场强强时自动降低增益、
弱时提高，
使检波输出幅度稳定；
同时设有 AFC（自动频率控制），
把本振频率锁向图像载频，
防止本振漂移导致图像对比下降或失色。
同步分离从全电视信号中取出行、
场同步脉冲，
经鉴相形成扫描电流去驱动偏转；
若同步不稳会出现图像滚动或撕裂，
AFC 与同步分离共同保证收、
发扫描严格同频同相。

彩色解码是模拟彩电的难点。
PAL 用梳状滤波器把色度中的 U、
V 两分量依行间相关性分离，
再经逐行倒相复原与解调得到色差信号；
当彩色副载波振荡器失锁或色同步信号受损时，
会出现“无彩色”或色调异常，
此时常设有“消色器（color killer）”自动关闭色度通道以免出现杂乱彩色噪声。
NTSC 因无色调逐行倒相补偿，
对本振频率精度与相位更敏感，
轻微频偏即引起色调偏移；
SECAM 则采用逐行轮换调频的色度，
抗微分相位失真较好。
伴音通道经鉴频得到音频，
与图像共用一个高频头但靠伴音内载波分离，
故调谐不准时常见“图像好、
伴音差”或反之。

模拟电视的显示核心多为阴极射线管（CRT）：
电子枪发射电子束，
经行/场偏转线圈扫描轰击荧光屏发光，
彩色管以红、
绿、
蓝三色荧光粉点或条按相加混色形成彩色。
三枪三束或单枪三束结构需精确“会聚”，
使三色光点重合；
偏转线圈决定几何形状，
聚焦极决定清晰度。
由于 CRT 依赖高压（数万伏）与磁场偏转，
其调试涉及聚焦、
会聚、
白平衡与几何校正，
是模拟彩电维修的重点。

常见问题多与上述环节相关：
散焦使图像发虚、
偏色源于白平衡或色通道失衡、
几何失真（桶形/枕形）来自偏转或枕校电路、
同步不良则因 AFC 或同步分离故障。
CRT 还需定期消磁（degauss）以消除地磁或外界磁场造成的色纯不良。
尽管 CRT 已被平板取代，
但其工作原理仍是理解模拟电视显示与维修的基础，
许多老设备仍在工业、
医疗与收藏领域使用。

\section{数字电视接收机}
\label{sec:digital_tv_receiver}

数字电视接收机接收数字电视信号，
具有以下优点：

\begin{itemize}
    \item \textbf{图像质量高}：无雪花、无重影
    \item \textbf{抗干扰能力强}：采用纠错编码
    \item \textbf{频道容量大}：数字压缩技术提高频谱利用率
    \item \textbf{功能丰富}：支持电子节目指南、交互式服务
\end{itemize}

数字电视标准包括DVB-T、
ATSC、
ISDB等。

数字电视接收机接收的是经过信源压缩与信道编码的调制载波。
在地面广播中，
欧洲 DVB-T/DVB-T2、
日本-巴西 ISDB-T、
美国 ATSC 以及中国 DTMB 是主要标准；
在有线电视中则大量采用 DVB-C（64-QAM 或 256-QAM）。
数字电视的核心优势来自“数字”二字：
信道中的噪声与干扰只要未造成误码超过纠错能力，
解码后图像质量就不会像模拟那样连续劣化，
因此“悬崖效应”（数字信号在门限之上几乎无损伤，
门下迅速崩溃）取代了模拟的渐进劣化。

数字地面电视普遍采用 OFDM（正交频分复用）以对抗多径。
OFDM 把高速数据流分散到大量彼此正交、
低速的子载波上，
配合循环前缀（guard interval）可有效消除码间干扰。
调制方式常随信道质量自适应在 QPSK、
16-QAM、
64-QAM 乃至 256-QAM 间切换，
子载波调制阶数越高频谱效率越高，
但对信噪比要求也越高。
DVB-T2 进一步引入更先进的 LDPC/BCH 纠错与旋转星座，
提升覆盖。
中国 DTMB 标准采用时域同步正交频分复用（TDS-OFDM），
以伪随机序列（PN）帧头同时完成信道估计与同步，
在单频网（SFN）中表现良好。

在数字域，子载波上的符号经解调后恢复为比特流，再经解交织、信道解码（如 RS、卷积码、LDPC、BCH 等）、解复用，得到视频与音频基本流，最后由信源解码器（如 MPEG-2、H.264/AVC、H.265/HEVC，或中国的 AVS/AVS+）还原为图像与声音。一个典型关系为：在 $M$ 进制 QAM 中，每符号携带 $\log_2 M$ 个比特，故在给定符号率 $R_s$ 下比特率为

\[
R_b=R_s\log_2 M
\]

机顶盒（STB）是数字电视时代的重要形态：
它内置调谐器、
信道解调与信源解码，
输出复合视频、
分量或 HDMI 信号给电视机，
并提供电子节目指南（EPG）、
条件接收（CA）、
交互点播等功能。
智能机顶盒还集成了操作系统与网络接口，
使电视机兼具互联网终端能力。

数字电视的核心是信源压缩与信道传输的协同。
视频经帧内/帧间预测、
变换与量化压缩为码流（如 MPEG-2、
H.264/AVC、
H.265/HEVC，
或中国的 AVS/AVS+），
多个节目的音视频基本流经打包与复用形成传输流（TS）；
信道编码再加入前向纠错（如 RS、
卷积码、
LDPC、
BCH）以抗误码。
地面广播采用 COFDM（编码正交频分复用），
把数据分散到大量正交子载波并加循环前缀以抗多径；
有线 DVB-C 则多采用多载波 QAM 在固定频道内高速传输。

单频网（SFN）是数字地面电视的重要特征：
多个发射台使用同一频率、
同一内容、
以 GPS 同源时钟严格同步发射，
依靠循环前缀吸收多径与邻台延时，
使覆盖区内信号叠加而非互扰，
从而高效利用频谱。
条件接收（CA）与电子节目指南（EPG）是运营侧功能：
CA 通过加扰与授权控制收视权限，
EPG 则在码流中携带节目时间表供用户浏览。
由于数字信号具有“门限效应”，
接收质量在门限之上几乎无损伤、
门限之下迅速崩溃，
故工程上常以误码率（BER）与调制误码率（MER）作为质量判据，
而非模拟的“雪花”程度。

\section{电视接收机电路原理}
\label{sec:tv_circuit_principle}

电视接收机的基本电路包括：

\begin{itemize}
    \item \textbf{高频头}：调谐、变频，输出固定中频
    \item \textbf{中频通道}：放大、滤波中频信号
    \item \textbf{视频处理}：解码视频信号，驱动显示
    \item \textbf{音频处理}：解码音频信号，驱动扬声器
    \item \textbf{同步电路}：分离行同步和场同步信号
\end{itemize}

彩色电视还包括色度解码电路，
还原红、
绿、
蓝三基色信号。

模拟与数字电视接收机在“前端—中频”部分结构相近，
差异主要体现在中频之后。
对模拟机而言，
中频通道是一组按残留边带特性设计的声表面波（SAW）滤波器与中放，
随后经视频检波得到全电视信号；
同步电路由行、
场同步分离与 AFC（自动频率控制）组成，
保证扫描与发射端严格同步，
否则会出现滚动或撕裂。
彩色解码部分（PAL 梳状滤波器与逐行倒相复原）已在上一节说明，
最终由矩阵电路把 Y、
U、
V 变换为驱动显像管的 R、
G、
B 三基色。

对数字机顶盒而言，
中频之后是 ADC 采样与全数字解调：
调谐器输出近基带的中频（或零中频 I/Q），
经 A/D 量化后由数字芯片完成 OFDM 解调、
均衡、
信道解码与信源解码，
最后经视频编码（如 CVBS、
YPbPr 或 HDMI）送显示。
相比模拟机，
数字机几乎把“图像质量”的责任从高频与中频模拟电路转移到数字解调算法与纠错码，
对模拟前端的线性度与相位噪声要求虽然仍高，
但对“选择性”的容限因数字滤波而宽松许多。

下表汇总模拟与数字电视接收机在若干关键方面的区别。

\begin{table}[htbp]
\centering
\caption{模拟电视与数字电视接收机主要区别}
\begin{tabular}{|p{4.2cm}|p{5.4cm}|p{5.4cm}|}
\hline
\textbf{比较项目} & \textbf{模拟电视接收机} & \textbf{数字电视接收机} \\
\hline
调制方式 & VSB-AM（残留边带）+ FM 伴音 & OFDM / QAM + 纠错编码 \\
\hline
中频之后处理 & 模拟检波、同步分离 & A/D 采样、数字解调解码 \\
\hline
抗干扰性 & 渐进劣化、怕多径重影 & 门限效应、强抗多径 \\
\hline
频道占用 & 每频道约 6\sim8MHz 传一套 & 同带宽可传多套（压缩） \\
\hline
典型标准 & PAL / NTSC / SECAM & DVB-T/T2、DTMB、ATSC、DVB-C \\
\hline
\end{tabular}
\end{table}

电视接收机的电源与高压部分是维修重点与安全风险所在。
开关电源把市电整流为低压供各电路，
行输出级再产生显像所需的高压与聚焦电压；
这部分元件承受大功率与高电压，
电容老化、
开关管击穿、
高压包匝间短路是常见故障，
检修必须断电放电、
使用隔离变压器并防止触电。
数字主板多为大规模 SoC，
集成调谐、
解调、
解码与接口，
故障常表现为黑屏、
无图、
卡顿或无法启动，
需结合指示灯、
软件日志与替换法定位。

模拟通道的故障则多集中在高频头、
中放、
色度解码与扫描电路，
可借助标准彩条信号、
示波器观测关键测试点逐级压缩。
现代一体机把模拟与数字处理高度集成，
板级维修多于元件级；
掌握“信号流向—测试点—典型波形/电压”的对应关系，
是高效排障的关键。
无论是 CRT 还是平板，
规范的安全操作与静电防护始终是第一位的。

\section{电视接收机调试与维护}
\label{sec:tv_receiver_maintenance}

调试项目：
\begin{itemize}
    \item \textbf{频率校准}：调整本振频率
    \item \textbf{图像调整}：亮度、对比度、色彩
    \item \textbf{伴音调整}：音量、音调
    \item \textbf{天线调整}：方向、增益
\end{itemize}

维护要点：
\begin{itemize}
    \item 定期清洁灰尘
    \item 检查电源电压
    \item 确保散热良好
    \item 避免长时间待机
\end{itemize}

调试电视接收机的首要环节是频率与增益校准。
模拟机需保证本振频率准确，
使图像中频落在 SAW 滤波器的标称位置；
本振偏移会导致图像对比度下降、
伴音失真甚至无彩色。
数字机顶盒虽由锁相环与数字同步自动锁定频率，
但在弱场强或强多径下仍需通过天线与放大器优化输入电平，
使 AGC 工作在正常区间（过高会引起非线性失真与互调，
过低则误码率上升）。

图像调整面向观看质量：
亮度（black level）、
对比度（white level）、
色饱和度与色调需依信号与显示器件设定；
对 PAL 机还需确认色度副载波振荡器锁定，
避免“无彩色”或“彩色爬行”（Hanover bars，
常因逐行倒相电路故障）。
伴音调整则检查鉴频零点、
音量音调与左右声道平衡。
天线调整对地面接收尤为关键：
方向、
极化与增益决定了输入场强与多径程度。

维护方面，
电视接收机（尤其含大功率电源与显示面板的整机）应定期清除内部积尘以防散热不良；
检查供电电压与保险，
避免电压浪涌损坏开关电源；
保证通风散热，
防止液晶背光或电源模块过热老化；
并尽量避免长期待机以减少待机功耗与元件电应力。
对使用多年的阴极射线管（CRT）机型，
还需注意高压器件与聚焦、
会聚的周期性检查。
规范调试与维护可显著延长设备寿命并维持图像与声音的接收质量。

除了上述项目，
现场调试还应注意天线馈线与接头质量：
劣质同轴电缆、
氧化或松动的 F 头会显著抬升噪声、
引入重影与误码，
往往比机内故障更常见。
室外天线需做好防雷与接地，
避免雷击损坏前端与人身安全事故。
对于数字接收，
应使用场强仪或机顶盒自带的信号质量（强度、
MER、
BER）指示判断“可解但临界”的状态，
单纯看信号强度不足以判断是否稳定收看。

维护中还应重视静电防护（ESD）与 firmware 升级：
现代数字机顶盒的调谐器与解码芯片对静电敏感，
插拔 HDMI 与天线前宜先放电；
厂商推送的软件升级可修复兼容性、
增加格式支持，
应定期执行但需确保供电稳定以防变砖。
安全方面，
含开关电源与高压组件的整机（尤其是旧式 CRT）拆机检修须先放电、
确认无高压后再操作。
综合来看，
电视接收机的可靠运行既靠出厂指标，
也靠规范的安装、
调试与日常维护。

调试与维修应配备基本仪器：
场强仪或带信号质量的机顶盒用于评估空中信号强度与质量，
标准信号发生器与彩条信号用于激励与比对，
示波器用于观测关键测试点波形与电压。
对数字接收，
重点看误码率（BER）、
调制误差比（MER）与星座图收敛度，
而非仅看信号强度条；
强度高但 MER 低往往表明多径或邻频干扰严重，
需调整天线方位、
极化或加滤波器。

常见故障存在一定“谱”：
无图像有噪声多为高频头或中频故障；
彩色丢失指向色度通道或副载波振荡；
图像不同步多为同步分离或 AFC；
数字机顶盒“有信号无节目”常为参数（频点、
符号率、
调制）设置错误或少载频。
建立按现象—可能部位—测试验证的排查流程，
结合图纸与经验库，
可显著提高维修效率并降低误换件率。
